\section{Why moving the query is a trade and propagating is not}
\label{sec:analysis}

\paragraph{A measured exchange rate.} Fix a rollout state $z_k$ and its anchor $y_0^{(k)}$.
Target construction moves $y_0^{(k)}$ a bounded distance and buys some reward. Define the
\emph{exchange rate} at that state as the realized local-reward gain per unit of relative
target displacement, $\widetilde R(\bar y_+)-\widetilde R(y_0)$ divided by
$\|\bar y_+-y_0\|_2/\|y_0\|_2$. Measured on $32$ held-out prompts with the trained
behavior policy, the reward-gradient construction gives
\begin{center}
\begin{tabular}{lcccc}
\toprule
$\sigma$ & $0.278$ & $0.465$ & $0.707$ & $0.913$\\
\midrule
exchange rate & $\mathbf{0.457}$ & $-0.46$ & $0.033$ & $0.0066$\\
\bottomrule
\end{tabular}
\end{center}
The query level chosen by \opsd is the only level where a unit of displacement buys a
substantial \emph{positive} amount: one step higher the construction is actively harmful,
and two or more steps higher it buys an order of magnitude less. \textbf{Moving the query away from it therefore buys trajectory
coverage at the price of target quality.} Propagation does not make that exchange: it
keeps the best construction site and reuses the resulting target elsewhere.

\paragraph{The trade is visible in the training objective itself.} If (A) optimizes a
worse target, it should reach a lower value of the training objective. It does. Over a
full $100$-update run, comparing each arm to its own matched canonical baseline at the
same rollout length,
\begin{center}
\begin{tabular}{lcc}
\toprule
& peak training reward & vs.\ matched canonical\\
\midrule
(A) query $\sigma\!:\,0.278\to0.465$ & $1.5964$ & $\mathbf{-0.0178}$\\
(C) \MSP & $1.6577$ & $\mathbf{+0.0415}$\\
\bottomrule
\end{tabular}
\end{center}
Moving the query lowers the height reached on the training objective; propagating raises
it. (Arms with a shortened rollout are excluded from this comparison because their
training reward is computed on a different sampling schedule and is not comparable. We use
this quantity only as a statement about the training objective, and draw no held-out
conclusion from it.)

\paragraph{Why propagation is not merely a smaller learning rate.} Equation~\eqref{eq:msploss}
divides the main-state term by $1+|\mathcal K|$, so one might suspect that \MSP simply
under-weights the original loss. It does not: halving the learning rate is significantly
\emph{harmful} in this setting ($-8.82$ MDE at $100$ updates), whereas \MSP is positive.
The extra states must therefore contribute useful gradient, and contribute enough to more
than offset the down-weighting of the main term.

\paragraph{Why the gain grows as the budget shrinks.} At $40$ steps both models traverse
every noise level with a small per-step $\Delta\sigma$, and the low-noise portion of the
trajectory can repair mid-noise errors. As the budget falls, each step spans a larger
$\Delta\sigma$ and the mid-noise transitions carry proportionally more of the generation;
this is exactly the region that propagation supervises and that the single-query baseline
never sees a loss at. Consistent with this, the gain is flat between $40$ and $10$ steps --
where neither model degrades -- and rises monotonically below $10$
(Fig.~\ref{fig:curves}).
